\chapter{Quem fez o Java}
\label{ap:pessoas}

O que as pessoas que construíram o Java escreveram e disseram: as
especificações, os design notes, os livros, as palestras e as entrevistas.
Cada seção tem uma bio de uma linha, com a página oficial de onde ela veio,
e uma lista curta do que ler ou assistir. Os links saem do \texttt{.bib}
(nenhum endereço foi copiado à mão); a lista completa está na
bibliografia, no apêndice~\ref{ap:bibliografia}. Quase tudo é gratuito e em
inglês.

\begin{description}
\item[Como foram escolhidos] Só entrou o que foi aberto e conferido em 18
  e 19/10/2026: título, autor e data na própria página. Os JEPs foram
  conferidos no snapshot do Internet Archive, porque o openjdk.org bloqueia
  acesso automatizado; o link impresso é o oficial.
\item[Quem tem pouco material] Algumas pessoas trabalharam mais em produto,
  compatibilidade ou comunidade do que escrevendo em público. Nesses casos
  a seção é curta e diz o que existe, sem enchimento.
\end{description}

\section{Comece pelo documentário}
\begin{itemize}
\obra{video:java-story}{Documentário de 2026, feito pela CultRepo, com a
  história do Java contada por quem o construiu. Segundo as reportagens da
  InfoWorld e do Foojay, aparecem Gosling, Bloch, Davidson, Lindholm,
  Polese, Schroer, Chander, Reinhold e Goetz: quase todo este capítulo.}
\obra*{artigo:infoworld-java-story}{Resenha com o resumo do filme: a
  ascensão, a queda e a volta do Java.}
\obra*{artigo:foojay-java-story}{Apresentação do filme pelo Foojay, com a
  lista de quem aparece e o papel de cada um.}
\obra{video:java30-livestream}{A transmissão oficial de seis horas dos 30
  anos do Java (2025), com trechos de Gosling, Reinhold, Goetz e Saab.}
\end{itemize}

\section{James Gosling}
Criador do Java: o Computer History Museum o fez Fellow em 2019 \enquote{pela
concepção, projeto e implementação da linguagem de programação Java}
(fonte: \citeobra{bio:gosling-chm}).

\begin{itemize}
\obra{paper:java-language-environment}{O white paper de 1996, com Henry
  McGilton: os objetivos de design (simples, orientada a objetos, robusta,
  portável) e o que foi deixado de fora do C e do C++. A melhor resposta a
  \enquote{por que o Java é assim?}.}
\obra{paper:feel-of-java}{Gosling, em primeira pessoa, sobre as decisões de
  design e por que o Java é uma linguagem \enquote{blue collar}.}
\obra{doc:jls25}{A especificação da linguagem, que ainda leva o nome dele
  entre os autores (com Joy, Steele, Bracha, Buckley, Smith e Bierman).}
\obra{livro:arnold2005}{O livro-texto da linguagem na época do J2SE~5.0
  (genéricos, enums, anotações), com Ken Arnold e David Holmes. Hoje vale
  como documento histórico.}
\obra{artigo:gosling-oral-history}{História oral do Computer History Museum
  (parte~1, transcrição gratuita): formação, CMU, Emacs e o caminho até o
  Green Project.}
\obra{video:gosling-foojay-30}{Podcast dos 30 anos do Java: das origens em
  controle de dispositivos aos servidores, e o que ele mudaria hoje.}
\end{itemize}

\section{Tim Lindholm}
Ex-Distinguished Engineer da Sun, contribuiu com a linguagem e foi o
arquiteto sênior da JVM; antes, trabalhou com VMs de Prolog (fonte:
\citeobra{bio:lindholm-informit}).

\begin{itemize}
\obra{doc:jvms25}{A especificação da JVM, que começou com ele e Frank
  Yellin: formato \texttt{.class}, verificação e cada instrução. É a
  referência dos capítulos de bytecode e de JVM da Parte~\ref{parte:java-jvm}.}
\obra{livro:lindholm2014}{A última edição impressa da JVMS (Java~SE~8).
  O item~1.1, \enquote{A Bit of History}, conta a origem da JVM.}
\obra{livro:lindholm1999}{A segunda edição, da era Sun. Serve para comparar
  com a JVM de hoje.}
\obra*{artigo:lindholm-register}{Entrevista de 2026: a reescrita às pressas
  das threads poucos dias antes da estreia do Java, e detalhes que o
  documentário não mostra.}
\end{itemize}

\section{Kim Polese}
Liderou o lançamento do Java na Sun e depois cofundou a Marimba (fonte:
\citeobra{bio:polese-chm}).

\begin{itemize}
\obra*{artigo:polese-talks-java}{Entrevista de 1996, no último dia dela na
  JavaSoft: por que a Sun \enquote{deu de graça} o Java, o licenciamento e
  os mercados que ela via para a linguagem.}
\item No documentário \citetitle{video:java-story}, como a primeira gerente
  de produto do Java.
\end{itemize}
Não achamos outro texto ou palestra dela sobre o Java que pudéssemos
conferir.

\section{James Duncan Davidson}
Na Sun, foi responsável pelas especificações Servlet~2.1/2.2 e JAXP~1.0; é
o autor original do Tomcat e do Ant (fonte: \citeobra{bio:davidson-ant}).

\begin{itemize}
\obra{doc:ant-faq}{O FAQ do Ant (texto do projeto, não dele) conta que o
  Ant nasceu dentro do Tomcat, para compilá-lo, e que o nome vem de
  \enquote{Another Neat Tool}, segundo Davidson. Veio antes do Maven e do
  Gradle.}
\item No documentário \citetitle{video:java-story}, sobre o Tomcat, o Ant e
  o Java no servidor.
\end{itemize}
Material técnico escrito por ele mesmo, aberto e conferido, não achamos.

\section{Carla Schroer}
Foi Director of Java Compatibility na Sun, segundo o Foojay (fonte:
\citeobra{artigo:foojay-java-story}). Não achamos página oficial dela.

\begin{itemize}
\item No documentário \citetitle{video:java-story}, sobre compatibilidade
  (o \enquote{write once, run anywhere}) e a disputa com a Microsoft.
\end{itemize}
Nenhuma palestra ou texto técnico dela foi encontrado e confirmado; o
documentário é, por enquanto, a fonte.

\section{Joshua Bloch}
Professor na CMU; foi chief Java architect no Google e distinguished
engineer na Sun, onde liderou o Collections Framework e as mudanças de
linguagem do JDK~5.0 (fonte: \citeobra{bio:bloch-informit}).

\begin{itemize}
\obra{livro:bloch2017}{Noventa itens de boas práticas: imutabilidade,
  \texttt{equals}/\texttt{hashCode}, genéricos, lambdas e streams. É o livro
  de Java mais citado nesta apostila.}
\obra{livro:goetz2006}{O clássico de concorrência, com Goetz, Peierls,
  Bowbeer, Holmes e Lea. Continua valendo com virtual threads.}
\obra{livro:bloch2005}{Quebra-cabeças curtos, com Neal Gafter, que expõem
  cantos estranhos da linguagem e da JVM. Bons para abrir uma aula.}
\obra{video:bloch-api-design}{Princípios de design de API, com exemplos bons
  e ruins tirados do próprio JDK.}
\obra{artigo:bumper-sticker-api}{As máximas da palestra em uma página, como
  \enquote{fácil de usar, difícil de usar errado}.}
\obra{artigo:binary-search-broken}{O overflow em \texttt{(low + high) / 2}
  que ficou anos no JDK: uma lição sobre invariantes e testes.}
\obra{video:bloch-effective-java-3e}{Os destaques da terceira edição do
  \emph{Effective Java}, com foco em lambdas e streams.}
\end{itemize}

\section{Brian Goetz}
Java Language Architect na Oracle; foi spec lead do JSR-335 (lambdas) e é o
autor de \emph{Java Concurrency in Practice} (fonte:
\citeobra{bio:goetz-inside}).

\subsection*{Lambdas e bytecode}
\begin{itemize}
\obra{doc:state-of-lambda}{Lambdas, interfaces funcionais, method
  references e default methods, explicados pelo designer. O documento não
  tem linha de autor: está hospedado no diretório pessoal dele em
  \texttt{cr.openjdk.org/\textasciitilde briangoetz}.}
\obra{doc:lambda-translation}{Como o \texttt{javac} traduz um lambda em
  \texttt{invokedynamic} e \texttt{LambdaMetafactory}. Leia com o
  \texttt{javap -c} aberto. Mesma hospedagem, também sem byline.}
\obra{video:goetz-lambda-hood}{A mesma tradução em palestra.}
\obra{doc:jep126}{O JEP que trouxe os lambdas ao Java~8 (Owner: Goetz).}
\obra{doc:jep484}{A API oficial para ler e gerar bytecode, que substitui o
  ASM dentro do JDK (autor: Goetz).}
\end{itemize}

\subsection*{Java moderno}
\begin{itemize}
\obra{artigo:dop-java}{Records, sealed types e pattern matching como
  modelagem de dados: tornar estados ilegais irrepresentáveis.}
\obra{doc:beyond-records}{Design note do Amber sobre levar a
  desestruturação a classes que não são records.}
\obra{doc:on-ramp}{Por que o Java simplificou o \enquote{hello world}.
  Ótimo para quem está começando.}
\obra{doc:jep286}{O \texttt{var}: motivação, limites e as alternativas
  descartadas (autor: Goetz).}
\obra{artigo:goetz-virtual-threads}{Virtual threads e o modelo
  thread-per-request.}
\obra{video:goetz-fp-vs-oo}{Funcional e orientado a objetos como
  ferramentas complementares.}
\obra{video:goetz-growing-java}{Como e por que a linguagem evolui (JVMLS
  2025).}
\end{itemize}

\subsection*{Valhalla e o futuro da JVM}
\begin{itemize}
\obra{doc:state-of-valhalla}{Value objects, layout de memória e por que
  objetos têm identidade.}
\obra{doc:defense-of-erasure}{Por que os genéricos do Java usam erasure,
  explicado por quem decide o futuro deles.}
\obra{video:goetz-valhalla-refactor}{Value classes, tipos null-restricted
  e strict initialization (Devoxx 2024).}
\end{itemize}

\section{Mark Reinhold}
Chief Architect do Java Platform Group na Oracle, onde trabalha no Java~SE e
no OpenJDK (fonte: \citeobra{bio:reinhold-inside}).

\begin{itemize}
\obra{doc:sotms}{O sistema de módulos explicado em uma leitura:
  \texttt{module-info}, readability, accessibility e módulos automáticos.}
\obra{video:reinhold-modules-one-lesson}{Os módulos ao vivo no terminal,
  em uma aula.}
\obra{doc:jep261}{A especificação de implementação do sistema de módulos
  (Owner: Reinhold).}
\obra{doc:jep403}{Por que o \texttt{--illegal-access} acabou e o que isso
  quebra em bibliotecas antigas que usam reflexão.}
\obra{video:reinhold-unsafe}{A história do \texttt{sun.misc.Unsafe} e por
  que ele vai sumindo.}
\obra{artigo:forward-faster}{A proposta de 2017 que criou a cadência de
  uma versão a cada seis meses.}
\obra{artigo:forward-even-faster}{A proposta de 2021 de uma LTS a cada dois
  anos.}
\obra{doc:jep3}{O processo de lançamento do JDK, por quem o desenhou.}
\obra{doc:leyden-beginnings}{Os objetivos do Project Leyden (startup,
  warmup, footprint): bom contraponto ao build time do Quarkus.}
\end{itemize}

\section{Paul Sandoz}
Architect no Java Platform Group da Oracle (fonte:
\citeobra{bio:sandoz-inside}). A página tem pouca bio; os JEPs e projetos
abaixo, que ele lidera, mostram o papel dele.

\begin{itemize}
\obra{doc:jep193}{VarHandles: acesso atômico e com ordenação de memória,
  a alternativa segura ao \texttt{Unsafe} (Owner: Sandoz; autor: Doug Lea).}
\obra{doc:jep338}{A Vector API: SIMD em Java puro, compilado pelo C2 para
  instruções vetoriais. Passou por várias rodadas de incubação, à espera do Valhalla.}
\obra{video:sandoz-vector-api}{Demonstração da Vector API e do código de
  máquina gerado, com Sandhya Viswanathan, da Intel.}
\obra{doc:babylon-code-models}{Code reflection: representar código Java
  simbolicamente para transformá-lo (GPU, SQL, diferenciação automática).}
\obra{video:sandoz-babylon}{O Project Babylon explicado pelo líder do
  projeto.}
\obra{video:sandoz-java-ai}{FFM, Vector API, Valhalla e Babylon aplicados a
  IA (JavaOne 2025).}
\obra*{doc:leyden-condensers}{Os \enquote{condensadores} do Leyden:
  transformar a aplicação antes de rodar, ideia parente do build time do
  Quarkus.}
\end{itemize}

\section{Georges Saab}
Senior Vice President do Java Platform Group na Oracle e chair do OpenJDK
Governing Board; na JavaSoft foi fundador do grupo do Swing e do
Java Web Start, e depois dirigiu a JRockit na BEA (fonte:
\citeobra{bio:saab-inside}).

\begin{itemize}
\obra{video:saab-rapid-innovation}{Por que o Java passou a lançar uma
  versão a cada seis meses (GOTO 2018).}
\obra*{video:saab-java-2024}{Keynote sobre o estado da plataforma em 2024.}
\obra*{video:saab-podcast-43}{Podcast: previsibilidade, que problemas o
  Java deve atacar e o foco recente em estudantes.}
\end{itemize}
Material escrito por ele é pouco; as palestras acima são o que há de mais
útil.

\section{Sharat Chander}
Hoje aparece como \enquote{Oracle Alum} no inside.java (fonte:
\citeobra{bio:chander-inside}); o Foojay o descreve como ex-Senior Director
de Java Product Management e Developer Engagement na Oracle.

\begin{itemize}
\obra{artigo:arrival-java21}{O anúncio oficial do JDK~21~LTS, com a lista
  de JEPs.}
\obra{artigo:java30-faceoff}{Retrospectiva dos 30 anos, com os recursos
  mais votados pela comunidade.}
\obra*{video:chander-25th}{Entrevista sobre a comunidade Java nos 25 anos.}
\end{itemize}

\section{Heather VanCura}
VP de Community Engagement na Oracle, à frente do programa Java Community
Process (JCP) (fonte: \citeobra{bio:vancura-qcon}).

\begin{itemize}
\obra{video:vancura-productivity}{Como adotar o Java atual e contribuir com
  a evolução da plataforma: JCP e grupos de usuários (JUGs).}
\obra*{artigo:vancura-qa}{Entrevista sobre o primeiro ano dela como chair do
  JCP e o papel das especificações.}
\obra{doc:jcp}{Onde ficam públicos os JSRs do Java~SE e os grupos de
  trabalho; qualquer pessoa pode acompanhar.}
\end{itemize}

\section{Para praticar no IDE: blog da JetBrains e JetBrains Guide}
\label{ap:jetbrains}

Não é gente que fez o Java, mas é onde muita gente aprende a usá-lo no dia
a dia: o blog do IntelliJ IDEA e os tutoriais gratuitos do JetBrains Guide.
São textos práticos, sempre com o IDE aberto, escritos por developer
advocates como Trisha Gee, Mala Gupta, Marit van Dijk, Helen Scott e Igor
Kulakov. Os mais úteis para cada capítulo também aparecem nas seções
\enquote{Onde aprofundar}.

% Todos os itens de jetbrains.bib entram nas listas por tema abaixo.
\nocite{artigo:jetbrains-writing-tests-junit5,
  artigo:jetbrains-migrating-junit4-junit5,
  artigo:jetbrains-testcontainers-spring-boot,
  artigo:jetbrains-ai-agent-flaky-tests,
  curso:jetbrains-guide-writing-junit5-tests,
  curso:jetbrains-guide-writing-tests-spock,
  curso:jetbrains-guide-testing-debugging,
  artigo:jetbrains-pattern-matching-5-examples,
  artigo:jetbrains-evolution-switch-java,
  artigo:jetbrains-project-loom-intellij,
  artigo:jetbrains-java-25-lts,
  artigo:jetbrains-do-you-really-know-java,
  artigo:jetbrains-java-still-relevant,
  artigo:jetbrains-thread-dumps-virtual-threads,
  artigo:jetbrains-debugging-java-intellij,
  artigo:jetbrains-sources-bytecode-debugging,
  artigo:jetbrains-println-debugging,
  artigo:jetbrains-debug-without-breakpoints,
  artigo:jetbrains-intellij-profiler,
  artigo:jetbrains-profiling-parrot-party,
  artigo:jetbrains-get-started-profiling,
  artigo:jetbrains-5-profiling-features,
  curso:jetbrains-guide-javax-jakarta,
  artigo:jetbrains-java-best-practices,
  artigo:jetbrains-code-smells-null,
  artigo:jetbrains-code-smells-mutation,
  artigo:jetbrains-code-smells-if,
  artigo:jetbrains-basic-refactorings,
  artigo:jetbrains-3-ways-refactor,
  curso:jetbrains-guide-refactoring}

\begin{description}
\item[Como foram conferidos] O blog da JetBrains bloqueia acesso
  automatizado; título, autor e data foram lidos no snapshot do Internet
  Archive, e o link impresso é o do blog. O JetBrains Guide foi aberto
  direto.
\item[O que ficou de fora] O post sobre String Templates: o recurso saiu
  em preview no JDK~21 e~22 e foi retirado (o JEP~465, terceira preview,
  está \emph{Closed/Withdrawn}); ensinar a sintaxe hoje seria ensinar algo
  que não existe no Java. Também saíram as pesquisas \emph{State of
  Developer Ecosystem}, com pouco conteúdo de Java, notas de versão do IDE e
  resumos mensais.
\end{description}

\subsection{Testes}
JUnit~5, migração do JUnit~4, Testcontainers e Spock; complementam a
Parte~\ref{parte:testes}.
\ondeachar*{jb-testes}

\subsection{Java moderno}
Pattern matching, \texttt{switch}, Project Loom e o Java~25; complementam
o capítulo de Java moderno.
\ondeachar*{jb-moderno}

\subsection{JVM, depuração e profiling}
Thread dumps com virtual threads, bytecode no depurador e o IntelliJ
Profiler; complementam os capítulos de JVM e de bytecode.
\ondeachar*{jb-jvm}

\subsection{Código limpo, refatoração e Jakarta}
A série \emph{Code Smells}, refatoração com o IDE e a migração de
\texttt{javax} para \texttt{jakarta}, que o Quarkus usa.
\ondeachar*{jb-limpo}

